\chapter{Principles and Schools of Yoga}
\label{ch:principles-schools}

\emph{How does yoga describe the human being, and what are the main paths it offers?}
The syllabus asks for two sets of ideas. The first is a group of models of the person: the three gunas, the fourfold inner instrument, the three bodies and the five sheaths, with the related lists of elements, vital airs, tissues and doshas. The second is the major schools or paths: jnana, bhakti, karma, Patanjala (raja) and hatha yoga, with mantra yoga beside them. These models are maps for understanding, not parts of the body; learn each as a short list with one plain word per item.

\begin{unverified}
This chapter is written from notes blocks whose fact check is still pending. Most points follow the official guidebook (its first, third and fifth chapters); where a list comes from the workshop slides or from standard Ayurveda rather than the guidebook, the text says so. A reviewer must check it before print.
\end{unverified}

\section{The three gunas (triguna)}

Samkhya teaches that nature, prakriti, is made of three qualities, the gunas:
\begin{itemize}
\item \textbf{Sattva}: purity, clarity, balance and harmony.
\item \textbf{Rajas}: activity, movement and restlessness.
\item \textbf{Tamas}: inertia, dullness, darkness and ignorance.
\end{itemize}
Every person and every thing is a mix of the three, and their balance shapes a person's mental and emotional state. The guidebook describes tamas and rajas as pulling a person down and sattva as lifting them up; the aim of practice is to grow sattva, for peace and spiritual growth. A sattvic person is focused, committed and positive; a rajasic person is active and restless; a tamasic person is lazy and dull.

\section{The inner instrument (antahkarana chatushtaya)}

The inner instrument has four parts:
\begin{itemize}
\item \textbf{Manas}: the mind that receives sense impressions and thoughts, and doubts.
\item \textbf{Buddhi}: the intellect that decides.
\item \textbf{Chitta}: the store of memory and past impressions (samskaras).
\item \textbf{Ahamkara}: the sense of ``I'', individuality.
\end{itemize}
This fourfold list follows the Upanishadic account given in the guidebook; Samkhya counts only three, without chitta. Bringing the four into harmony supports clarity and self-awareness.

\section{The three bodies and the five sheaths}

\subsection*{Tri-sharira}
The person is described as three bodies, one inside the other:
\begin{itemize}
\item \textbf{Sthula sharira}: the gross, physical body.
\item \textbf{Sukshma sharira}: the subtle body of vital energy and mind.
\item \textbf{Karana sharira}: the causal body.
\end{itemize}

\subsection*{Panchakosha}
The five-sheath model goes back to the Taittiriya Upanishad and was set out as the five koshas by Adi Shankaracharya. The sheaths, from the outside in, are:
\begin{itemize}
\item \textbf{Annamaya kosha}: the food sheath, the physical body.
\item \textbf{Pranamaya kosha}: the sheath of vital energy, the bridge between body and mind.
\item \textbf{Manomaya kosha}: the sheath of mind, thinking and emotion.
\item \textbf{Vijnanamaya kosha}: the sheath of higher understanding and intuition.
\item \textbf{Anandamaya kosha}: the bliss sheath, closest to the Self but still a sheath; the Self lies behind all five.
\end{itemize}

The guidebook joins the two models. The gross body is the annamaya sheath; the subtle body is made of the pranamaya, manomaya and vijnanamaya sheaths; the causal body is the anandamaya sheath. So the subtle body is more than the mind: it includes the vital sheath.

The model also shows how the practices work, from the outside in: asana acts on the food sheath, pranayama on the energy sheath, and meditation on the sheaths of mind and understanding. The guidebook uses this model in its account of yoga for health, and the five-sheath theory is the one usually named as the basis of yoga in health care.

\ycbfig{ycb1-fig-kosha}

\section{Related lists}

\subsection*{The five great elements (pancha mahabhuta)}
Space (\emph{ākāśa}), air (\emph{vāyu}), fire (\emph{agni}), water (\emph{jala}) and earth (\emph{pṛthvī}), each arising from the one before it in this order. In the body, space gives the cavities; air gives movement, pulsation and conduction; fire gives metabolism; water gives binding and the joints; earth gives the organs, mass and firmness.

\subsection*{Subtle elements and senses}
\textit{Source: workshop slides.} Each element has a subtle element (\emph{tanmātra}) and a sense organ: sound and the ears for space; touch and the skin for air; form or light and the eyes for fire; taste and the tongue for water; smell and the nose for earth. The five organs of action (\emph{karmendriya}) are speech, hands, feet, excretion and reproduction.

\subsection*{The seven tissues (sapta dhatu)}
The guidebook says the food sheath is built of seven dhatus. The names come from Ayurveda: rasa (nutrient fluid), rakta (blood), mamsa (muscle), meda (fat), asthi (bone), majja (marrow) and shukra (reproductive tissue).
\ycbfig{ycb1-fig-dhatus}

\subsection*{The three doshas (tridosha)}
\textit{Source: workshop slides.} \emph{Dosha} means that which goes out of balance. Vata is dry, light, cold, rough and unsteady; pitta is hot, sharp, light, liquid and quick; kapha is unctuous, cold, heavy, sluggish, soft and steady.

\subsection*{The five vital airs (pancha prana)}
\begin{itemize}
\item \textbf{Prana}: the upper region, from the head to the navel.
\item \textbf{Apana}: from the navel down to the base of the trunk.
\item \textbf{Udana}: rising from the navel to the head.
\item \textbf{Samana}: the region of the navel.
\item \textbf{Vyana}: spread through the whole body.
\end{itemize}
The guidebook also names five minor airs (\emph{upaprāṇa}): naga (belching and hiccup), kurma (blinking), krikala (hunger, thirst, sneezing and coughing), devadatta (yawning and drowsiness) and dhananjaya (said to remain after death).
\ycbfig{ycb1-fig-pranas}

\subsection*{Chakras and other quick facts}
Two chakra locations are often asked: ajna lies between the eyebrows and vishuddhi at the base of the throat. The Mandukya Upanishad links the waking, dream and deep-sleep states with the sounds A, U and M. The four aims of life (\emph{puruṣārtha}) are dharma, artha, kama and moksha; samadhi is not one of them. In Vedantic language the Self is \emph{sat-cit-ānanda}, being, consciousness and bliss; avidya, ignorance of one's true nature, is the first of Patanjali's five afflictions; and Ishvara in the Yoga Sutra is a special purusha.
\ycbfig{ycb1-fig-chakras}

\section{The major schools of yoga}

The guidebook names six prominent paths: jnana, bhakti, karma, raja, hatha and mantra yoga. A person usually blends them. The common protocol also groups the paths by what they use: body (karma), mind (jnana), emotion (bhakti) and energy (kriya).

\subsection*{Jnana yoga, the path of knowledge}
The path of discerning the real from the unreal, non-dual (Advaita Vedanta) and contemplative. It has three stages: \emph{śravaṇa}, listening to the teaching; \emph{manana}, reflection and reasoning; and \emph{nididhyāsana}, repeated meditation that makes the conviction firm. The student prepares through the fourfold discipline (\emph{sādhana catuṣṭaya}) and studies Vedanta under a realised teacher. The standard fourfold list is discrimination (viveka), dispassion (vairagya), the six virtues (shama, dama, uparati, titiksha, shraddha, samadhana) and the longing for freedom (mumukshutva).

\subsection*{Bhakti yoga, the path of devotion}
The path of intense, unconditional love for the divine. The Narada Bhakti Sutra is a principal text, and kirtan, devotional chanting, is perhaps its most important technique. The Gita (seventh chapter, verse 16) names four kinds of devotee: the distressed, the inquisitive, the seeker of wealth and the wise. The Bhagavata Purana lists nine forms of devotion: hearing, chanting, remembering (\emph{smaraṇa}, constant remembrance), serving, worship, salutation, service as a servant, friendship and self-surrender.

\subsection*{Karma yoga, the path of action}
The path of working without attachment to the result. The guidebook draws four principles from the Gita: act from a sense of duty; stay free of strong attachment to the outcome; do not let worry about results disturb the mind while working; and accept success and failure with evenness. Offering the fruit of action to God melts the ego into selfless action.

\subsection*{Raja yoga, Patanjali's path}
The royal yoga, based on the eight limbs of the Yoga Sutra, aims at stilling the movements of the mind so that the true self stands in its own nature. It works on the mind. The Hatha Pradipika calls hatha yoga a staircase leading to raja yoga, so raja yoga is also used for the final state rather than a separate technique.

\subsection*{Hatha yoga}
The path that works through the body and breath, with cleansing actions, asanas, breath control and locks, as a base for meditation. The guidebook traces it to Tantra, and notes that hatha systems usually begin with the cleansing actions while Patanjali's limbs begin with yama and niyama; this is a difference in where practice starts, not a rejection of yama and niyama. Its texts include the Hatha Pradipika of Svatmarama, the Gheranda Samhita, the Goraksha Samhita of Gorakshanath and the Shiva Samhita. The Pradipika explains \emph{ha} as the sun and \emph{tha} as the moon, the union of the pingala and ida channels. Reading hatha as ``stubbornness'' is a mistake.

\subsection*{Mantra yoga}
The yoga of sound. The word \emph{mantra} is explained as that which protects (\emph{tra}) the mind (\emph{man}). A mantra is a sacred sound used for deep meditation; om is called the great mantra.

\subsection*{Everyday examples}
Each path can be shown with one example: studying and reflecting (jnana); prayer or song (bhakti); service done without expecting thanks (karma); a daily sitting practice (raja); an asana and breath routine (hatha).

\begin{recall}
\begin{itemize}
\item Gunas: sattva, rajas, tamas; practice grows sattva.
\item Inner instrument: manas, buddhi, chitta, ahamkara.
\item Bodies: sthula, sukshma, karana. Sheaths: annamaya, pranamaya, manomaya, vijnanamaya, anandamaya.
\item Elements in order: space, air, fire, water, earth. Pranas: prana, apana, udana, samana, vyana.
\item Paths: jnana, bhakti, karma, raja, hatha, mantra.
\item Jnana stages: shravana, manana, nididhyasana.
\end{itemize}
\end{recall}

\subsection*{Worked examples}
\begin{itemize}
\item \textbf{Model answer: explain the five sheaths and which practice works on each.} Annamaya, the physical body, through asana; pranamaya, the energy sheath, through pranayama; manomaya and vijnanamaya, mind and understanding, through meditation; anandamaya, the bliss sheath, innermost, with the Self behind all five.
\item \textbf{A student asks which path is ``the right one''.} Explain that the paths suit different temperaments and that most people blend them; a daily routine of asana and breath, some study, and service already touches several.
\end{itemize}

\subsection*{Check yourself}
\begin{enumerate}
\item Name the three gunas.
\item Name the four parts of the inner instrument.
\item Which sheath does pranayama work on?
\item Which sheaths make up the subtle body, according to the guidebook?
\item Name the five great elements in order.
\item Which dhatu is blood?
\item What does the word dosha mean?
\item Which prana governs the region of the navel?
\item Which chakra lies between the eyebrows?
\item Name the three stages of jnana yoga.
\item Which path uses action without attachment to its fruit?
\item What do ha and tha stand for in hatha yoga?
\end{enumerate}

\subsection*{Answers}
\begin{enumerate}
\item Sattva, rajas, tamas.
\item Manas, buddhi, chitta, ahamkara.
\item The pranamaya kosha, the energy sheath.
\item The pranamaya, manomaya and vijnanamaya sheaths.
\item Space, air, fire, water, earth.
\item Rakta.
\item That which goes out of balance.
\item Samana.
\item Ajna.
\item Shravana, manana and nididhyasana.
\item Karma yoga.
\item The sun (ha) and the moon (tha).
\end{enumerate}

\paragraph{Sources.} The official guidebook for Level I (2016, older scheme), first, third and fifth chapters; the Common Yoga Protocol, Ministry of AYUSH; the Bhagavad Gita and the Bhagavata Purana (devotees and forms of devotion); standard Ayurvedic usage for the names of the dhatus; the workshop slides for the tanmatras and doshas (to be checked); the Board syllabus for Level 1.
